\section*{Abstract}
\addcontentsline{toc}{section}{Abstract}
Federated recommendation separates shared item learning from local user personalization,
but useful recommendations can deteriorate when whole-user differential privacy (DP)
requires clipping and noisy aggregation. This project studies whether a smaller shared
representation improves that trade-off and, more broadly, how injected noise becomes
disturbance to item matrices, personalized scores, and rankings. A reproducible
MovieLens-100K pipeline compares centralized Bayesian Personalized Ranking (BPR), a
user-weighted centralized control, and federated full-matrix and low-rank models with
and without DP. Five-seed historical comparisons find low-rank private models better
than the private full-matrix reference in four of sixteen multiplicity-adjusted
comparisons; all four nonprivate low-rank variants underperform the converged full
model. These results motivate a separately frozen, equal-budget comparison of full,
jointly trained two-factor, and fixed-public-factor parameterizations.

The fixed-factor extension obtains no positive adjusted primary comparison and loses
eight of sixteen comparisons to joint factor training. A user-level private popularity
baseline exceeds every collaborative point estimate at target $\epsilon\leq4$ in
that extension. Conditional moment calculations and measured perturbations show why
factor-coordinate noise norms are inadequate proxies for recommendation damage:
effective noise depends on factor geometry, local user scale, clipping, and pairwise
margins. Pulse-and-reset studies initially reveal large two-factor trajectory
differences, but an orthogonal alignment control substantially reduces them. Thirty
transferred MovieLens-1M training runs and 200 further probes reproduce the dependence
on the way noise is coupled across rebalanced factor trajectories. A separate screen
of 3,200 item-pair records finds a maximum absolute sign-flip probability error of
$7.67\times10^{-6}$ for a variance-matched Gaussian approximation and rejects
non-Gaussian tails as a practically consequential explanation in the sampled states.

The completed contribution is an audited experimental study and a more precise
interpretation of noise diagnostics in personalized low-rank recommendation. It
does not establish a new state-of-the-art algorithm or a confirmed claim of
paper-level novelty. The report includes the full methodological progression,
unfavorable findings, privacy scope, statistical qualifications, mathematical
derivations, and a separately generated supplementary results document.

\medskip
\noindent\textbf{Keywords:} federated recommendation; user-level differential privacy;
matrix factorization; BPR; effective noise; personalization; factor alignment;
reproducible evaluation.

\subsection*{Scope of this submission}
The main report explains the completed research and includes essential tables and
all nine curated figures. The companion \emph{Supplementary Results and Experimental
Record} provides complete retained utility tables, primary and secondary contrasts,
diagnostic summaries, and source mappings. Numerical results come from frozen saved
outputs; preparation of this report did not retrain models or score additional test
targets. MovieLens-100K is a repeatedly used development benchmark. MovieLens-1M
validation has been evaluated; its test targets remain unscored.
